% !PRINTED 191
\weberpage{191}
%<<CONT
and we must therefore set the determinant $\alpha\delta - \beta\gamma = \pm 1$.
If one computes the determinant $|A'|$, one obtains for it
$(\alpha\delta - \beta\gamma)^3$, so that, if we take account only
of the proper orthogonal substitutions, we must set
\weq{13}{\alpha\delta - \beta\gamma = 1}
while
\weq{14}{\alpha\delta - \beta\gamma = -1}
corresponds to the improper orthogonal substitutions.
By this the substitution $A'$ is completely determined, and is
reduced to the binary substitution
\weq{15}{(\eta_1, \eta_2) = \begin{pmatrix} \alpha, & \beta \\ \gamma, & \delta \end{pmatrix} (\xi_1, \xi_2),\quad \alpha\delta - \beta\gamma = \pm 1,}
or, if $\eta_1 : \eta_2 = \eta$, $\xi_1 : \xi_2 = \xi$ is set,
to the linear fractional substitution
\weq{16}{\eta = \dfrac{\alpha\xi + \beta}{\gamma\xi + \delta}.}
The following point remains to be noted. The two substitutions
$A, A'$ are transformations of one another, and thus correspond
to each other one-to-one.

The numbers $\alpha, \beta, \gamma, \delta$ are determined by $A'$
only up to a common sign. If, then, we consider the
linear substitution
\weq{17}{\begin{pmatrix} \alpha, & \beta \\ \gamma, & \delta \end{pmatrix}}
with one of the two conditions (13) and (14), then the substitution
$A'$, and hence also $A$, is indeed uniquely determined by this;
but conversely to each $A'$ there correspond two
substitutions of the form (17):
\[
\begin{pmatrix} \alpha, & \beta \\ \gamma, & \delta \end{pmatrix}
\quad\text{and}\quad
\begin{pmatrix} -\alpha, & -\beta \\ -\gamma, & -\delta \end{pmatrix}.
\]
But we again obtain a one-to-one correspondence if we
combine these two substitutions into a common concept, which
we denote by $A''$.

Finally, as regards the substitution (16), which we denote by
\weq{18}{A''' = \begin{pmatrix} \alpha, & \beta \\ \gamma, & \delta \end{pmatrix}}
and, following \S.\,38, conceive as a substitution of the ratios,
we can remove a common factor of
$\alpha, \beta, \gamma, \delta$
%CONT>>
